\chapter{Laparoscopy}

\section*{Overview}
Laparoscopy is a minimally invasive procedure in which a camera and surgical instruments are passed through small incisions, called ports, in the abdomen. It allows the surgeon to see inside the abdomen and pelvis without a large incision. It can be used for diagnosis, treatment, or both.

\section*{Alternative Names}
Keyhole surgery, minimally invasive abdominal surgery, diagnostic laparoscopy, therapeutic laparoscopy

\section*{Principle}
A narrow telescope with a light and camera is inserted through a small incision. Carbon dioxide gas is used to inflate the abdomen so that organs can be seen clearly. Instruments are passed through additional ports to inspect and treat structures.

\section*{History}
The first laparoscopic procedures were performed in the early 1900s, and technical advances expanded the technique through the 1970s. The first laparoscopic cholecystectomy in 1985 transformed abdominal surgery. Laparoscopy is now the standard approach for many operations.

\section*{Why It Is Done}
Diagnostic laparoscopy is used to inspect the pelvic and abdominal organs when imaging is unclear. Therapeutic laparoscopy treats conditions such as gallbladder disease, hernias, endometriosis, and ovarian cysts. It is often preferred because recovery is faster than with open surgery.

\section*{Is This a Primary or Secondary Test or Procedure}
It is a primary procedure for many elective operations and a secondary tool when noninvasive tests are inconclusive. It may be diagnostic, therapeutic, or both in the same operation.

\section*{How It Is Performed}
The patient is placed under general anesthesia, and the abdomen is inflated with carbon dioxide. A port is created at the navel for the camera, and additional small ports are placed for instruments. The surgeon inspects or treats the organs and then removes the ports and closes the small wounds.

\section*{Preparation}
Fasting for six to eight hours before the procedure and emptying the bladder and bowels as directed. Blood-thinning medications are stopped beforehand. A pregnancy test is done when appropriate.

\section*{Contraindications and Precautions}
\begin{itemize}
  \item Severe cardiopulmonary disease that cannot tolerate gas inflation
  \item Extensive prior abdominal surgery with adhesions
  \item Uncorrected coagulopathy
  \item Late pregnancy
\end{itemize}

\section*{Risks}
Bleeding, infection, and injury to abdominal organs or blood vessels are possible. Port-site hernias can develop at the incision sites. Gas-related shoulder and abdominal discomfort is common and temporary.

\section*{Aftercare and Recovery}
Most patients go home the same day or stay overnight. Mild pain, bloating, and shoulder ache from the gas are expected. Activity can usually be resumed within a few days to a week.

\section*{Results and Interpretation}
A diagnostic laparoscopy confirms the presence or absence of disease seen at surgery. A therapeutic laparoscopy is successful when the condition is corrected and the patient recovers without complications.

\section*{Expected Results}
Normal-appearing abdominal and pelvic organs with no findings requiring treatment. For therapeutic procedures, successful completion without conversion to open surgery or major complications.

\section*{Follow-up}
A follow-up visit is scheduled to review the surgical findings and check the wounds. Heavy lifting is restricted for several weeks to prevent port-site hernias. Additional testing is guided by the findings and any tissue sent to the laboratory.

\section*{References}
\begin{enumerate}
  \item MedlinePlus. Laparoscopy. U.S. National Library of Medicine; 2025.
  \item Current Medical Diagnosis and Treatment. McGraw-Hill; 2025.
\end{enumerate}

\clearpage
